\section{The fourth digamma moment and harmonic Laurent data}
\label{bhr:quartic-harmonic:sec}

The incoming \emph{Cubic Tornheim Derivatives and Harmonic Stieltjes Constants}
report, Section~7,
Question~6 \cite{bhr:Cubic}, explicitly asks for a power-to-harmonic formula at the fourth
power. The following theorem answers that question. It is separate from
the quartic \emph{directional} Tornheim layer, which has already been
worked out in later incoming reports. No claim of worldwide novelty or
finite reduction to ordinary zeta constants is intended.

Write $H_n^{(r)}=\sum_{j=1}^n j^{-r}$, $H_n=H_n^{(1)}$, and
\[
 H_n(1,1)=\sum_{n\ge j>k\ge1}\frac1{jk}
         =\frac{H_n^2-H_n^{(2)}}2.
\]
The strict harmonic Dirichlet series used here are
\[
 Z_2(s,1)=\sum_{n=1}^{\infty}\frac{H_{n-1}}{n^s},\qquad
 Z_3(s,1,1)=\sum_{n=1}^{\infty}\frac{H_{n-1}(1,1)}{n^s},
 \qquad \Re s>1.
\]
They have the local expansions
\begin{align}
 Z_2(1+u,1)&=u^{-2}+\gamma u^{-1}
       +\frac{\gamma^2-\zeta(2)}2+\eta u+O(u^2),
       \label{bhr:quartic-harmonic:eta}\\
 Z_3(1+u,1,1)&=u^{-3}+\gamma u^{-2}
       +\frac{\gamma^2-\zeta(2)}{2u}
       +\frac{\gamma^3-3\gamma\zeta(2)+2\zeta(3)}6
       +\beta_1u+O(u^2).
       \label{bhr:quartic-harmonic:beta}
\end{align}
In particular $\beta_1$ is a completely specified strict depth-three
harmonic Laurent coefficient; it is not a coefficient of the inclusive
harmonic series. The constant term is also fixed by the Mellin lemma below.
The pole terms in the second expansion follow directly
from $H_{n-1}=\log n+\gamma+O(n^{-1})$ and
$H_{n-1}^{(2)}=\zeta(2)+O(n^{-1})$, by subtraction of zeta derivatives.
The convention for the ordinary Stieltjes constants is
\[
 \zeta(1+u)=u^{-1}+\sum_{j\ge0}\frac{(-1)^j\gamma_j}{j!}u^j.
\]

\begin{theorem}[A quartic power-to-harmonic identity]
\label{bhr:quartic-harmonic:main}
Use the unit-coordinate Hadamard finite part at the origin and put
$Q_4=\operatorname{FP}\int_0^1\psi(x)^4\,dx$. Then
\begin{align}
 Q_4={}&18\zeta(4)+12\zeta(3)-12\gamma\zeta(2)+4\zeta'(3)
       +12(\gamma^2-\zeta(2))\gamma_1
       +24\gamma\eta-24\beta_1.
       \label{bhr:quartic-harmonic:main-eq}
\end{align}
An equivalent absolutely convergent arithmetic formula is
\begin{align}
 Q_4={}&18\zeta(4)+12\zeta(3)-12\gamma\zeta(2)+4\zeta'(3)
          +12\gamma_3-24\zeta(2)\gamma_1\notag\\
 &+12\sum_{n=1}^{\infty}
       \frac{\bigl(\psi(n)^2+\psi'(n)-\log^2 n\bigr)\log n}{n}.
       \label{bhr:quartic-harmonic:arithmetic}
\end{align}
If $Q_3=\operatorname{FP}\int_0^1\psi(x)^3\,dx$, the already proved
cubic bridge \cite{bhr:Cubic} $Q_3=3\zeta(2)-6\eta-6\gamma\gamma_1$ eliminates $\eta$:
\begin{equation}
 Q_4+4\gamma Q_3
   =18\zeta(4)+12\zeta(3)+4\zeta'(3)
       -12(\gamma^2+\zeta(2))\gamma_1-24\beta_1.
       \label{bhr:quartic-harmonic:combined}
\end{equation}
\end{theorem}

\subsection{A normalized Mittag--Leffler expansion}

For $n\ge0$ set
\begin{align*}
 b_n&=H_n-\gamma,& c_n&=\zeta(2)+H_n^{(2)},&
 e_n&=H_n^{(3)}-\zeta(3),\\
 a_n&=c_n-b_n^2,& d_n&=6b_n^2-4c_n,&
 r_n&=-4b_n^3+12b_nc_n-4e_n.
\end{align*}
Let
\[
 C_* =\gamma^4-12\gamma^2\zeta(2)+12\gamma\zeta(3)+11\zeta(4),
 \qquad
 P_n(z)=\frac1{(n+z)^4}-\frac{4b_n}{(n+z)^3}
          +\frac{d_n}{(n+z)^2}+\frac{r_n}{n+z}.
\]

\begin{lemma}[Quartic principal parts with no unknown entire term]
\label{bhr:quartic-harmonic:ML}
For $z\notin\{0,-1,-2,\ldots\}$,
\begin{equation}
 \psi(z)^4=P_0(z)+C_*+
     \sum_{n=1}^{\infty}\bigl(P_n(z)-P_n(0)\bigr).
     \label{bhr:quartic-harmonic:ML-eq}
\end{equation}
The series converges normally on every compact set avoiding the poles.
\end{lemma}

\begin{proof}
The digamma recurrence and its convergent expansion at zero \cite{bhr:DLMFG} give
\[
 \psi(-n+\varepsilon)
   =-\varepsilon^{-1}+b_n+c_n\varepsilon+e_n\varepsilon^2
       +(\zeta(4)+H_n^{(4)})\varepsilon^3+O(\varepsilon^4).
\]
Raising this expression to the fourth power gives exactly $P_n$ as
its principal part. The constant coefficient at $n=0$ is
\[
 \gamma^4-12\gamma^2\zeta(2)+6\zeta(2)^2
                +12\gamma\zeta(3)-4\zeta(4)=C_*.
\]
This establishes the poles but does not on its own establish the entire
remainder. To determine that remainder, let $f(z)=\psi(z)^4-P_0(z)$.
For fixed $z$, apply the residue theorem to
$z f(w)/(w(w-z))$ on $|w|=N+\tfrac12$.
The standard digamma asymptotic in the right half-plane, and
$\psi(w)=\psi(1-w)-\pi\cot(\pi w)$ in the left half-plane, give
$f(w)=O(\log^4 N)$ uniformly on these circles. The cotangent is
uniformly bounded there: away from the real axis this follows from its
exponential formula, while near the real axis these half-integer circles
remain a fixed positive distance from the integer poles. Thus the
contour integral is $O(\log^4N/N)$ and tends to zero. The residues are
$f(z)$, $-f(0)$, and $-(P_n(z)-P_n(0))$, proving the formula.
Finally $b_n=O(\log n)$, $d_n=O(\log^2n)$ and
$r_n=O(\log^3n)$, and the simple-pole difference is $O(n^{-2})$
on each fixed compact set. These estimates prove normal convergence.
\end{proof}

\subsection{The exact summation-by-parts cancellation}

For $m\ge2$ and $n\ge1$,
\[
 \int_0^1\bigl((n+x)^{-m}-n^{-m}\bigr)\,dx
     =\frac{n^{1-m}-(n+1)^{1-m}}{m-1}-n^{-m}.
\]
Integrating \eqref{bhr:quartic-harmonic:ML-eq}, and retaining the finite
parts of its three higher-order poles at zero, yields
\begin{equation}
 Q_4=K+\sum_{n=1}^{\infty}r_n
      \left(\log\left(1+\frac1n\right)-\frac1n\right),
      \label{bhr:quartic-harmonic:K-eq}
\end{equation}
where
\begin{equation}
 K=\gamma^4-18\gamma^2\zeta(2)+32\gamma\zeta(3)
       +\frac{13}{2}\zeta(4)+12\zeta(3)-12\gamma\zeta(2).
       \label{bhr:quartic-harmonic:K}
\end{equation}
For clarity, the ordinary Euler sums needed in this simplification are
\begin{equation}
 \sum_{n\ge1}\frac{H_n}{n^2}=2\zeta(3),\qquad
 \sum_{n\ge1}\frac{H_n}{n^3}=\frac54\zeta(4),\qquad
 \sum_{n\ge1}\frac{H_n^{(2)}}{n^2}=\frac74\zeta(4),\qquad
 \sum_{n\ge1}\frac{H_n^2}{n^2}=\frac{17}4\zeta(4).
 \label{bhr:quartic-harmonic:euler}
\end{equation}
These are classical ordinary convergent sums. One short derivation of
the three weight-four formulas is as follows. The shuffle and stuffle
products of $\zeta(2)^2$ give respectively
$2\zeta(2,2)+4\zeta(3,1)$ and $2\zeta(2,2)+\zeta(4)$;
hence $\zeta(3,1)=\zeta(4)/4$ and
$\zeta(2,2)=3\zeta(4)/4$.
The iterated-integral substitution $t\mapsto1-t$ gives
$\zeta(2,1,1)=\zeta(4)$. Expanding the three harmonic numerators
in strict sums now gives exactly the displayed values. The first,
weight-three formula follows from the same shuffle/stuffle method or
from the elementary symmetric integral
$\int_0^1\log x\log(1-x)/(x(1-x))\,dx=2\zeta(3)$.

The decisive finite difference is
\begin{equation}
 r_n-r_{n-1}=\frac{12a_{n-1}}n+\frac4{n^3}.
 \label{bhr:quartic-harmonic:residue-difference}
\end{equation}
It follows by substituting
$b_n=b_{n-1}+n^{-1}$, $c_n=c_{n-1}+n^{-2}$,
and $e_n=e_{n-1}+n^{-3}$; no limiting argument is involved.
Thus finite summation by parts gives
\begin{align}
 \sum_{n=1}^N r_n\log\left(1+\frac1n\right)
 &=r_N\log(N+1)-12\sum_{n=1}^N\frac{a_{n-1}\log n}{n}
                      -4\sum_{n=1}^N\frac{\log n}{n^3}.
 \label{bhr:quartic-harmonic:SBP}
\end{align}

We also need the constant term in $\sum_{n=1}^N r_n/n$.
Here is an exact finite identity that makes this normalization
checkable. Put $X=H_N$, $Y=H_N^{(2)}$, $Z=H_N^{(3)}$,
$W=H_N^{(4)}$, and
\[
 S_{12}=\sum_{n=1}^N\frac{H_n}{n^2},\quad
 S_{13}=\sum_{n=1}^N\frac{H_n}{n^3},\quad
 S_{22}=\sum_{n=1}^N\frac{H_n^{(2)}}{n^2},\quad
 S_{112}=\sum_{n=1}^N\frac{H_n^2}{n^2}.
\]
Then
\begin{align}
 \sum_{n=1}^N\frac{r_n}{n}
 ={}&-X^4+4\gamma X^3+(-6\gamma^2+6\zeta(2))X^2
       +(4\gamma^3-12\gamma\zeta(2)+4\zeta(3))X\notag\\
 &+6X^2Y-12\gamma XY-4XZ
       +(-6\gamma^2+6\zeta(2))Y-16\gamma Z-11W\notag\\
 &+24\gamma S_{12}-12S_{112}+20S_{13}+6S_{22}.
 \label{bhr:quartic-harmonic:finite-sum}
\end{align}
For example this identity follows by taking finite differences of
$H_n^4$, $H_n^3$, $H_n^2H_n^{(2)}$, $H_nH_n^{(2)}$,
and $H_nH_n^{(3)}$, then collecting terms. Both sides vanish at
$N=0$, and their exact finite differences are $r_N/N$.

Let $L=\log N$. Using \eqref{bhr:quartic-harmonic:euler} in
\eqref{bhr:quartic-harmonic:finite-sum}, together with
$H_N=L+\gamma+O(N^{-1})$, gives
\begin{equation}
 \sum_{n=1}^N\frac{r_n}{n}
 =-L^4+12\zeta(2)L^2+C+o(1),\qquad
 C=\gamma^4-18\gamma^2\zeta(2)+32\gamma\zeta(3)
       -\frac{23}2\zeta(4).
 \label{bhr:quartic-harmonic:C}
\end{equation}
Also
$r_N\log(N+1)=-4L^4+24\zeta(2)L^2+o(1)$.
Since $a_{n-1}=2\zeta(2)-\log^2n+O(\log n/n)$, the limit
\begin{equation}
 A=\lim_{N\to\infty}\left\{
  \sum_{n=1}^N\frac{a_{n-1}\log n}{n}
      -\zeta(2)L^2+\frac{L^4}{4}\right\}
 \label{bhr:quartic-harmonic:A}
\end{equation}
exists. Equations \eqref{bhr:quartic-harmonic:K-eq},
\eqref{bhr:quartic-harmonic:SBP}, and \eqref{bhr:quartic-harmonic:C}
now imply
\begin{equation}
 Q_4=18\zeta(4)+12\zeta(3)-12\gamma\zeta(2)
                         +4\zeta'(3)-12A.
 \label{bhr:quartic-harmonic:before-Laurent}
\end{equation}
Every divergent polynomial has been retained until this cancellation.

\subsection{Convergent arithmetic and the Laurent coefficient}

At a positive integer,
\[
 a_{n-1}=2\zeta(2)-\psi(n)^2-\psi'(n).
\]
Moreover $\psi(n)^2+\psi'(n)-\log^2n=O(\log n/n)$.
Using the standard Stieltjes defining limits in
\eqref{bhr:quartic-harmonic:A} gives
\[
 A=2\zeta(2)\gamma_1-\gamma_3
   -\sum_{n=1}^{\infty}
          \frac{\bigl(\psi(n)^2+\psi'(n)-\log^2n\bigr)\log n}{n}.
\]
The final series is absolutely convergent because its summand is
$O(\log^2n/n^2)$. Substitution proves
\eqref{bhr:quartic-harmonic:arithmetic}.

To identify the harmonic coordinate, introduce
\[
 F(s)=\sum_{n=1}^{\infty}
    \frac{\psi(n)^2+\psi'(n)-\log^2n}{n^s},\qquad\Re s>0.
\]
This series and every fixed spectral derivative converge normally
on compact subsets of the stated half-plane. Expanding the harmonic
numerator in strict sums gives, initially for $\Re s>1$,
\begin{equation}
 F(s)=2Z_3(s,1,1)-2\gamma Z_2(s,1)
                   +(\gamma^2+\zeta(2))\zeta(s)-\zeta''(s).
 \label{bhr:quartic-harmonic:deformation}
\end{equation}
This also supplies the required local continuation. The coefficient
of $u$ at $s=1+u$ is
\[
 F'(1)=2\beta_1-2\gamma\eta
                 -(\gamma^2+\zeta(2))\gamma_1+\gamma_3.
\]
On the other hand it is the negative of the convergent sum in
\eqref{bhr:quartic-harmonic:arithmetic}. Inserting this relation proves
\eqref{bhr:quartic-harmonic:main-eq}. Combining it with the cubic bridge
proves \eqref{bhr:quartic-harmonic:combined}.

\subsection{A Mellin normalization for every height-one coefficient}

The following classical Mellin mechanism (compare the harmonic-zeta framework in \cite{bhr:Kargin}) supplies an independent
realization of the coefficient used above. Its all-depth formulation
also states precisely where new global information first enters.
Let $d\ge1$, and write
\[
 Z_d(s,\{1\}^{d-1})
   =\sum_{n=1}^{\infty}\frac{H_{n-1}(\{1\}^{d-1})}{n^s},
 \qquad
 \frac1{\Gamma(1+u)}=\sum_{j\ge0}c_ju^j.
\]
An empty harmonic word has value one. For $t>0$ put
\[
 B(t)=-\log(1-e^{-t}),\qquad
 V_d(t)=B(t)^d-\mathbf1_{0<t<1}(-\log t)^d.
\]
Since $V_d(t)=O(t(1+|\log t|^{d-1}))$ at zero and decays
exponentially at infinity, every logarithmic moment of $V_d(t)/t$
converges absolutely.

\begin{lemma}[All height-one finite parts and regular coefficients]
\label{bhr:quartic-harmonic:mellin}
For $|u|<1$, as an identity of meromorphic functions,
\begin{equation}
 \Gamma(1+u)Z_d(1+u,\{1\}^{d-1})
   =\frac1{u^d}+\frac{u}{d!}
       \int_0^\infty t^{u-1}V_d(t)\,dt.
 \label{bhr:quartic-harmonic:mellin-generator}
\end{equation}
Consequently the principal part and constant coefficient together are
$\sum_{j=0}^{d}c_j u^{j-d}$. For every $k\ge1$, the exact regular
coefficient is
\begin{equation}
 [u^k]Z_d(1+u,\{1\}^{d-1})
  =c_{d+k}+\frac1{d!}\sum_{j=1}^k
       \frac{c_{k-j}}{(j-1)!}
       \int_0^\infty\frac{V_d(t)\log^{j-1}t}{t}\,dt.
 \label{bhr:quartic-harmonic:mellin-coefficients}
\end{equation}
In particular the finite part is $c_d$, while the coefficient of $u$
is $c_{d+1}+d!^{-1}\int_0^\infty V_d(t)\,dt/t$.
\end{lemma}

\begin{proof}
The elementary multiple-polylogarithm identity
$\operatorname{Li}_{\{1\}^{d-1}}(z)
=(-\log(1-z))^{d-1}/(d-1)!$ gives, first for $\Re s>1$,
\[
 \Gamma(s)Z_d(s,\{1\}^{d-1})
  =\int_0^\infty t^{s-1}
       \frac{B(t)^{d-1}}{(d-1)!(e^t-1)}\,dt.
\]
Subtract $(-\log t)^{d-1}/((d-1)!t)$ on $(0,1)$.
Its Mellin integral at $s=1+u$ is $u^{-d}$. The subtracted
kernel is $-V_d'(t)/d!$; no jump term occurs at $t=1$, since
$(-\log1)^d=0$. Integrating that remainder by parts produces
exactly the integral in \eqref{bhr:quartic-harmonic:mellin-generator}.
The boundary terms vanish for $\Re u>-1$, because of the bounds on
$V_d$. Normal convergence of that integral and all its spectral
derivatives proves continuation and coefficient extraction.
\end{proof}

For the present depth-three coordinate this yields the particularly
simple convergent formula
\begin{equation}
 \beta_1=
 \frac{\gamma^4}{24}-\frac{\gamma^2\zeta(2)}4
       +\frac{\gamma\zeta(3)}3+\frac{\zeta(4)}{16}
       +\frac16\int_0^\infty\frac{V_3(t)}t\,dt.
 \label{bhr:quartic-harmonic:beta-mellin}
\end{equation}
Likewise $\eta=c_3+\tfrac12\int_0^\infty V_2(t)\,dt/t$.
These formulas evaluate the harmonic coefficients through elementary
convergent kernels; they do not identify those kernels with finite
combinations of ordinary constants. They also confirm the constant
coefficient displayed in \eqref{bhr:quartic-harmonic:beta} directly.

\subsection{A normalized primitive and every differentiated identity}

Put
\[
 d_0=6\gamma^2-4\zeta(2),\qquad
 r_0=4\gamma^3-12\gamma\zeta(2)+4\zeta(3).
\]
For $\Re z>0$ with the principal logarithm, define
\begin{align}
 \mathcal P_4(z)={}&-\frac1{3z^3}-\frac{2\gamma}{z^2}
                      -\frac{d_0}{z}+r_0\log z+C_*z\notag\\
 &+\sum_{n=1}^{\infty}\Bigg[
   \frac1{3n^3}-\frac1{3(n+z)^3}-\frac z{n^4}\notag\\
 &\hspace{11mm}-4b_n\left\{\frac1{2n^2}-\frac1{2(n+z)^2}
                                      -\frac z{n^3}\right\}\notag\\
 &\hspace{11mm}+d_n\left\{\frac1n-\frac1{n+z}-\frac z{n^2}\right\}
       +r_n\left\{\log\left(1+\frac zn\right)-\frac zn\right\}
                                      \Bigg].
 \label{bhr:quartic-harmonic:primitive}
\end{align}
The normally convergent braces are the integrals, vanishing at zero,
of the subtracted principal parts. Consequently
\[
 \mathcal P_4'(z)=\psi(z)^4,\qquad
 \operatorname{FP}\int_0^b\psi(x)^4\,dx=\mathcal P_4(b)
      \quad(b>0),\qquad \mathcal P_4(1)=Q_4.
\]
The constant term of $\mathcal P_4$ at the singular origin is zero;
this fixes the integration constant without an unspecified branch term.

For every integer $m\ge1$, differentiation of
\eqref{bhr:quartic-harmonic:ML-eq} gives the further exact identity
\begin{align}
 &\sum_{m_1+m_2+m_3+m_4=m}
   \frac{m!}{m_1!m_2!m_3!m_4!}
      \prod_{j=1}^4\psi^{(m_j)}(z)\notag\\
 &\quad=(-1)^m m!\sum_{n=0}^{\infty}
 \left\{
  \frac{\binom{m+3}{3}}{(n+z)^{m+4}}
  -\frac{4b_n\binom{m+2}{2}}{(n+z)^{m+3}}
  +\frac{(m+1)d_n}{(n+z)^{m+2}}
  +\frac{r_n}{(n+z)^{m+1}}
 \right\}.
 \label{bhr:quartic-harmonic:derivatives}
\end{align}
The convergence is normal off the digamma poles, since the slowest
term is $O(\log^3n/n^{m+1})$. At $z=1$, the left side is a finite
polynomial in $\gamma$ and ordinary zeta values.

\paragraph{Independent checks.}
The exact checking program verifies the principal parts, the residue
finite difference, the full finite identity
\eqref{bhr:quartic-harmonic:finite-sum}, and the cancellations producing
\eqref{bhr:quartic-harmonic:main-eq}. A separate numerical calculation
integrates a locally Taylor-subtracted digamma fourth power and compares
it with the convergent arithmetic formula, whose tail is evaluated by
Bernoulli asymptotics and Hurwitz-zeta derivatives. At working precision
80 digits, the arithmetic-tail evaluation differs from quadrature by
less than $3.4\cdot10^{-62}$. Independent Mellin evaluations of $\eta$
and $\beta_1$, inserted into the main identity, differ from the same
quadrature by less than $2.1\cdot10^{-74}$. The common value begins
\[
 Q_4=13.5101203760158274116062407128304484419327985333019186009860964\ldots.
\]
These are numerical diagnostics, not certified enclosures or substitutes
for the proof.

\paragraph{Next precise targets.}
The theorem determines $\beta_1$ from the cubic and quartic finite parts,
but does not express it in a prescribed algebra of ordinary constants.
For the fifth power, the residue polynomial has degree four in the
finite harmonic data, and its first difference should produce both
height-one depth-four Laurent data and additional harmonic numerators.
The correct next task is to derive that exact difference and identify the
minimal \emph{stated} family of Laurent coefficients. A second direction
is a common Hurwitz shift: the shifted strict harmonic numerator and the
integration endpoint must be normalized together, since replacing
$n$ by $n+a$ only in the outer denominator does not shift the harmonic
prefix or the finite-part coordinate.
